\chapter{Documentação de programas de CLP}
\label{cap:documentacao}

\begin{objetivos}
Ao final deste capítulo você deve ser capaz de:
\begin{itemize}
  \item listar os artefatos mínimos que compõem a documentação de um projeto;
  \item montar uma tabela de alocação de E/S;
  \item documentar um bloco de programa com cabeçalho padronizado;
  \item definir uma rotina de versionamento e backup compatível com a realidade da planta;
  \item reconhecer a documentação como requisito contratual de entrega.
\end{itemize}
\end{objetivos}

\section{Por que documentar}

Um programa sem documentação só funciona enquanto seu autor está disponível. Na prática industrial,
isso significa: parada mais longa na manutenção corretiva, dificuldade de formar novo operador,
incapacidade de provar o que foi alterado e risco em auditoria de segurança.

A própria apostila de origem já registrava, no capítulo de programação em Ladder, que ``uma prática
indispensável é a elaboração das tabelas de alocação ou mapa de memória dos dispositivos de
entrada/saída''. Esta apostila retoma e amplia essa prática, porque ela é o item que costuma faltar
justamente nas instalações mais antigas --- as que mais precisam de manutenção.

\index{documentação de programa}\index{tabela de alocação}

\section{Artefatos mínimos}

\begin{table}[H]
  \centering
  \caption{Documentação mínima de um projeto de CLP.}
  \label{tab:artefatos-documentacao}
  \begin{tabular}{p{4.2cm} p{5.6cm} p{4.6cm}}
    \toprule
    Artefato & Conteúdo & Quem usa \\
    \midrule
    Arquitetura de automação
      & Diagrama dos controladores, redes, remotas e estações
      & Manutenção, elétrica, novos projetos \\
    Tabela de alocação de E/S
      & TAG, descrição, localização e endereço de cada ponto
      & Manutenção, operação, programação \\
    Lista de tags e variáveis
      & Nome, tipo, escopo, retenção e descrição de cada variável
      & Programação e manutenção \\
    Descritivo operacional
      & Como o processo funciona e o que acontece em cada condição
      & Operação e programação \\
    Diagrama de sequência ou de fluxo
      & Passos do processo e condições de transição
      & Operação, manutenção e testes \\
    Listagem comentada
      & Programa com comentários de bloco e de linha
      & Manutenção \\
    Manual de operação e de alarmes
      & Modos, comandos, alarmes e o que fazer em cada um
      & Operação \\
    Rotina de backup e versionamento
      & Frequência, local, nomenclatura e responsável
      & Manutenção e gestão \\
    \bottomrule
  \end{tabular}
\end{table}

\section{Tabela de alocação de E/S}

É o documento mais consultado na manutenção, porque liga três mundos: o instrumento em campo, o ponto
físico no cartão e o nome que aparece no programa.

\begin{table}[H]
  \centering
  \caption{Modelo de tabela de alocação de E/S (exemplo).}
  \label{tab:tabela-alocacao}
  \begin{tabular}{l l l l}
    \toprule
    TAG & Descrição & Localização & Endereço \\
    \midrule
    PSL-100 & Chave de pressão baixa & Topo do tanque pressurizado & \texttt{DI1} \\
    TT-400  & Transmissor de temperatura & Saída do misturador & \texttt{AI1} \\
    FS-200  & Chave de fluxo & Saída de óleo do aquecedor & \texttt{DI2} \\
    TV-400  & Válvula de temperatura & Ao lado do tanque & \texttt{DO1} \\
    \bottomrule
  \end{tabular}
\end{table}

Note a estrutura: \textbf{TAG} (identificação de instrumentação), \textbf{descrição} (o que é),
\textbf{localização} (onde está fisicamente) e \textbf{endereço} (onde o programa o encontra). Sem a
localização, o endereço não serve para quem está em campo com o painel aberto; sem o endereço, o TAG
não serve para quem está programando.

\begin{atencao}
Mantenha a tabela de alocação atualizada \emph{junto} com o programa. Em instalação em que a tabela é
feita uma vez, no comissionamento, e nunca revisada, o técnico do plantão descobre o endereço
testando ponto por ponto com a máquina em operação --- exatamente o cenário que as boas práticas do
Capítulo~\ref{cap:boas-praticas} querem evitar.
\end{atencao}

\section{Documentação dentro do programa}

Documentar não é só produzir documento externo: é deixar o próprio programa legível. Um cabeçalho
padronizado por POU resolve a maior parte do problema.

\begin{lstlisting}[language=, caption={Modelo de cabeçalho de POU.}, label={lst:cabecalho-pou}]
(* ------------------------------------------------------------------
   POU:        FB_Controle_Bomba
   Tipo:       Bloco de função
   Objetivo:   Comandar bomba com protecao de nivel e registro de horas
   Entradas:   Liga (BOOL), Desliga (BOOL), Nivel (REAL, m), NivelMin (REAL, m)
   Saidas:     Bomba (BOOL), Horas (DINT, h), Falha (BOOL)
   Regras:     Liga somente se Nivel >= NivelMin. Nivel < NivelMin
               desliga a bomba. Falha bloqueia liga ate reset.
   Retencao:   Horas (RETAIN)
   Conforme:   IEC 61131-3, 4a ed. (2025) - blocos padronizados
   Autor/Data: <nome> / <data>      Revisao: 02 (2026-09-22)
   ------------------------------------------------------------------ *)
\end{lstlisting}

Três regras práticas para comentário:

\begin{enumerate}
  \item \textbf{comente o porquê, não o quê}: ``liga a bomba'' o código já diz; ``protege contra
        cavitação'' diz o que ninguém deduz;
  \item \textbf{numere as revisões e registre o que mudou} --- a pergunta ``o que mudou depois do
        comissionamento?'' é a mais frequente na manutenção;
  \item \textbf{descreva cada tag na lista de variáveis}, não apenas no programa.
\end{enumerate}

\section{Versionamento e backup}

\begin{table}[H]
  \centering
  \caption{Rotina mínima de backup e versionamento.}
  \label{tab:backup}
  \begin{tabular}{p{4.4cm} p{10.2cm}}
    \toprule
    Item & Recomendação \\
    \midrule
    O que versionar
      & O arquivo-fonte do projeto (que abre no ambiente de programação), não apenas o PDF ou o
        arquivo compilado. Sem a fonte, não há manutenção possível \\
    Frequência
      & Sempre após alteração testada, antes e depois de qualquer intervenção em campo e em
        periodicidade fixa (por exemplo, mensal) \\
    Nomenclatura
      & Nome da máquina, número de revisão, data e responsável: \texttt{ENV01\_rev07\_2026-09-22\_WSV} \\
    Onde guardar
      & Fora do controlador, em pelo menos dois locais, um deles fora da planta. Backup só no
        computador da sala de controle não é backup \\
    Quem responde
      & Um responsável nomeado. Backup sem responsável não acontece \\
    Restauração
      & Testada pelo menos uma vez por ano. Backup nunca testado é suposição, não garantia \\
    \bottomrule
  \end{tabular}
\end{table}

\paragraph{Alteração em campo.} Programação \emph{on-line} (Capítulo~\ref{cap:funcionamento}) é o
ponto em que a documentação se perde com mais facilidade: altera-se a lógica com a máquina rodando e
ninguém registra. A regra prática: \textbf{toda alteração em campo gera uma revisão nova do arquivo,
com data e autor}, mesmo que seja de uma linha.

\section{Documentação como entrega contratual}

A documentação discutida aqui é a mesma que a Seção de especificação e compra exige do fornecedor
(Capítulo~\ref{cap:porte}). Ela deve estar prevista no contrato, com lista de artefatos, formato e
prazo --- porque no fim da obra a pressa é grande e o que não está no contrato não é entregue.

Para a instituição de ensino, vale a mesma lógica: trabalho de conclusão, projeto integrador e
programa didático entregues sem documentação ensinam ao aluno o hábito errado.

\begin{pratica}
Monte a documentação de um trabalho prático seu, com quatro entregas:
(a) arquitetura de automação em diagrama;
(b) tabela de alocação de E/S no modelo da Tabela~\ref{tab:tabela-alocacao};
(c) cabeçalho de POU para cada bloco do programa, no modelo da
listagem~\ref{lst:cabecalho-pou};
(d) rotina de backup preenchida na Tabela~\ref{tab:backup}, com nome do responsável.
Depois troque o material com um colega e peça que ele explique o que o seu programa faz sem abrir o
código. O que ele não conseguir explicar é o que falta documentar.
\end{pratica}

\begin{exercicios}
\begin{enumerate}[label=\alph*)]
  \item Explique por que a tabela de alocação de E/S é o documento mais consultado na manutenção.
  \item Uma instalação tem backup semanal do programa, guardado no computador da sala de controle.
        Aponte dois problemas dessa rotina e como corrigi-los.
  \item Por que versionar apenas o PDF do programa não é suficiente?
  \item Escreva o cabeçalho de POU para um bloco que controla uma esteira com sensor de presença e
        contador de peças.
  \item Dê um exemplo de comentário que explica ``o quê'' e outro que explica ``porquê''. Qual dos
        dois tem mais valor e por quê?
  \item Uma alteração de uma linha foi feita em campo, com a máquina rodando, sem registro. Descreva
        o risco que isso cria para o próximo mantenedor.
  \item Liste o que deve constar em contrato de fornecimento de sistema de automação sobre
        documentação.
  \item Por que restaurar um backup de teste, uma vez por ano, é diferente de apenas fazer backup?
\end{enumerate}
\end{exercicios}
